% 多圈重整化、复合算符混合、算符乘积展开与有限温度边界条件
单圈原始发散图通常只有一个整体紫外发散，需要相应的局域反项吸收。进入多圈以后，同一图内部还会出现嵌套或重叠的发散真子图。若等整张图积分完成后再一次性减去极点，子图发散会被带入更高阶图中，无法由局域参数重定义逐阶吸收。局域性因此要求一种递归减除：先对所有真子图完成重整化，再处理缩约图剩余的整体发散。

对一张单粒子不可约图 $\Gamma$，记维数正则化后的裸振幅为 $I_\Gamma$。令线性算符 $K$ 只抽取所选减除方案中的局域紫外发散部分；例如在修正的最小减除方案（$\overline{\rm MS}$）中，$K$ 抽取所有 $1/\bar\epsilon^n$ 极点。对没有发散真子图的原始发散图，反项定义为
\begin{equation*}
 C(\Gamma)=-K I_\Gamma.
\end{equation*}
存在发散真子图时，先把所有真子图的局域反项插回原图，定义不完全 $R$ 运算
\begin{equation}
 \boxed{
 \bar R I_\Gamma
 =I_\Gamma
 +\sum_{\{\gamma_i\}}
 \left[\prod_i C(\gamma_i)\right]
 I_{\Gamma/\{\gamma_i\}}.}
 \label{eq:rbar-operation-dp68}
\end{equation}
其中求和遍历两两不重叠的紫外发散真子图集合，$\Gamma/\{\gamma_i\}$ 表示把每个 $\gamma_i$ 缩成一个局域顶角。完整减除的递归结构可表述如下。
\begin{proposition}[局域反项的递归减除]
固定正则化与线性紫外减除投影 $K$，并假设真子图递归减除后剩余的整体发散是 $K$ 所抽取的局域部分；红外奇异性已被另行调节。以两两不重叠的发散真子图集合定义 $\bar R$，其中各 $C(\gamma)$ 按较小子图递归定义。则整体反项和剩余振幅满足
\begin{equation}
 \boxed{
 C(\Gamma)=-K\bar RI_\Gamma,
 \qquad
 RI_\Gamma=(1-K)\bar RI_\Gamma}
 \label{eq:r-operation-dp11}
\end{equation}
\end{proposition}
\begin{proof}
先说明递归的代数结构。对没有发散真子图的图，$\bar RI_\gamma=I_\gamma$，于是 $C(\gamma)=-KI_\gamma$。假设所有较小真子图的反项已经定义，$\bar RI_\Gamma$ 是原图加上所有两两不重叠真子图集合的反项插入；把剩余整体反项定义为 $-K\bar RI_\Gamma$，立即有
\[
RI_\Gamma=\bar RI_\Gamma+C(\Gamma)=(1-K)\bar RI_\Gamma.
\]
图的真子图严格较小，所以归纳在有限步内结束。局域性与红外调节是假设；该代数恒等式本身不会证明这些分析条件。

逐层展开 $C(\gamma)=-K\bar RI_\gamma$ 时，顶层只选互不重叠的子图，下一层只选其内部的真子图。所得子图集合因而两两不重叠或互相包含；反过来，任一这样的森林由其最大真子图及各自的内部森林唯一确定。这给出递归项与兼容森林的一一对应，嵌套减除必须从内到外执行。重叠而互不包含的两个子图不能同时出现，分别属于不同森林。完整 $R$ 的最后一次整体减除还允许森林包含 $\Gamma$ 本身；只取真子图的森林描述的是 $\bar R$。

对只有一个发散真子图的两圈图，上述归纳的第一步可显式计算如下。

设两圈振幅为 $I_\Gamma$，其中 $\gamma$ 是一个对数发散的单圈真子图。对子图定义

\begin{equation*}
 C(\gamma)=-K I_\gamma.
\end{equation*}
在只有这一个发散真子图的情形，正文\eqnrefs{eq:rbar-operation-dp68} 化为
\begin{align*}
 \bar R I_\Gamma
 &=I_\Gamma+C(\gamma)I_{\Gamma/\gamma}\\
 &=I_\Gamma-(K I_\gamma)I_{\Gamma/\gamma}.
\end{align*}
第二项在每一个外层动量处抵消子图的大动量局域极点，因此 $\bar R I_\Gamma$ 只可能还保留整张图共同放大时的整体局域发散。正文\eqnrefs{eq:r-operation-dp11} 已把这一剩余整体发散的局域反项写成 $C(\Gamma)=-K\bar R I_\Gamma$。将该反项插回后便有
\begin{align*}
 R I_\Gamma
 &=\bar R I_\Gamma+C(\Gamma)\\
 &=(1-K)\bar R I_\Gamma,
\end{align*}
即正文\eqnrefs{eq:r-operation-dp11}。多个互不重叠或彼此嵌套的真子图只是在\eqnrefs{eq:rbar-operation-dp68} 中增加允许的子图集合，随后同样只对剩余整体发散作最后一次局域减除。
\end{proof}
第一式只抽取已经消除子图发散后剩余的整体局域极点；在上述分析假设下，第二式留下有限的重整化图。

嵌套与重叠子图可以用 Zimmermann 森林公式统一组织。用于完整 $R$ 的森林 $F$ 是一组彼此不重叠或互相包含的发散子图；若 $\Gamma$ 整体发散，也允许 $\Gamma$ 本身属于森林。以下明确取 BPHZ 的局域 Taylor 减除算符 $T_\gamma$，其取值阶数为相应表面发散度；它与前面最小减除中的极点投影 $K$ 是不同的方案实现。
\begin{corollary}[递归局域减除的森林组织]
假设 BPHZ 子图减除满足前述局域性条件，且所论外动量及红外调节允许相应 Taylor 减除。森林中的嵌套算符按从内到外的顺序作用，两两不重叠的减除可独立执行。则有限递归展开为
\begin{equation}
 \boxed{
 RI_\Gamma
 =\sum_{F\in\operatorname{Forests}(\Gamma)}
 \left(\prod_{\gamma\in F}(-T_\gamma)\right)I_\Gamma.}
 \label{eq:zimmermann-forest-dp11}
\end{equation}
\end{corollary}
\begin{proof}
对无发散真子图的图，森林仅含空集以及整体发散时的 $\{\Gamma\}$，求和给 $I_\Gamma-T_\Gamma I_\Gamma$。归纳到一般图：任一真子图森林的最大元素两两不重叠；每个最大元素内的子森林由较小图的归纳式给出。因此递归地选择最大元素、再选择其内部森林，恰好列举所有兼容的真子图森林，每项出现一次。最后整体减除把每个真森林项分别乘以 $I-T_\Gamma$，产生不含及包含 $\Gamma$ 的两项，得到所述完整森林和。互相重叠而不包含的子图不可能属于同一最大元素集合，因此不会被重复同时减除。此证明建立有限减除式的组合组织；紫外局域性与积分收敛仍依赖陈述中列明的分析假设。
\end{proof}
空森林给出原图；含一个子图的森林给一次减除；嵌套森林按从内到外的顺序作用；含整体 $\Gamma$ 的项完成最后的整体减除。定理的关键假设是紫外发散部分在外动量和轻质量上可作局域 Taylor 展开，因此 $T_\gamma$ 抽取的是有限阶局域多项式。这一局域性正是多圈反项仍能写回原拉格朗日量允许算符的原因。

重整化局域算符还需要处理重合点奇异性，因此场与耦合的重整化之后必须进一步引入复合算符混合。

分离时空点的 Green 函数完成重整化后，局域复合算符仍会产生额外的重合点奇异性。若一组算符 $\mathcal O_i$ 具有相同守恒量子数，重整化一般不是逐个乘一个常数，而是矩阵混合：
\begin{equation}
 \boxed{\mathcal O_{i,B}=Z_{ij}^{\mathcal O}\mathcal O_{j,R}.}
 \label{eq:composite-operator-mixing-dp11}
\end{equation}
定义反常维数矩阵
\begin{equation*}
 \gamma^{\mathcal O}
 \equiv (Z^{\mathcal O})^{-1}\mu\frac{\dd Z^{\mathcal O}}{\dd\mu},
\end{equation*}
即可得到
\begin{equation*}
 \mu\frac{\dd}{\dd\mu}\mathcal O_{i,R}
 =-\gamma_{ij}^{\mathcal O}\mathcal O_{j,R}.
\end{equation*}
若有效拉格朗日量含 $C_i(\mu)\mathcal O_i(\mu)$，物理作用量不能依赖任意重整化尺度，因此 Wilson 系数沿相反方向演化：
\begin{equation}
 \boxed{
 \mu\frac{\dd C_i}{\dd\mu}
 =C_j\gamma_{ji}^{\mathcal O}.}
 \label{eq:wilson-rg-dp11}
\end{equation}

\begin{proposition}[算符混合的对偶系数演化]
设重整化算符基封闭于线性混合，系数与算符的收缩表示同一个尺度不变的相互作用。则有以下结论（原式 \eqnrefs{eq:wilson-rg-dp11}）：
\begin{equation*}
 \boxed{
 \mu\frac{\dd C_i}{\dd\mu}
 =C_j\gamma_{ji}^{\mathcal O}.}
\end{equation*}
\end{proposition}
\begin{proof}
由

\begin{equation*}
 \mathcal O_B=Z^{\mathcal O}\mathcal O_R
\end{equation*}
并固定裸理论参数，
\begin{equation*}
 0=\mu\frac{\dd Z^{\mathcal O}}{\dd\mu}\mathcal O_R
 +Z^{\mathcal O}\mu\frac{\dd\mathcal O_R}{\dd\mu}.
\end{equation*}
左乘 $(Z^{\mathcal O})^{-1}$，并定义
\begin{equation*}
 \gamma^{\mathcal O}
 =(Z^{\mathcal O})^{-1}\mu\frac{\dd Z^{\mathcal O}}{\dd\mu},
\end{equation*}
得到
\begin{equation*}
 \mu\frac{\dd\mathcal O_{i,R}}{\dd\mu}
 =-\gamma_{ij}^{\mathcal O}\mathcal O_{j,R}.
\end{equation*}
有效作用量中的同一局域项写成 $C_i\mathcal O_i$。物理量不依赖任意重整化尺度，所以
\begin{align*}
 0
 &=\mu\frac{\dd}{\dd\mu}(C_i\mathcal O_i)\\
 &=\left(\mu\frac{\dd C_i}{\dd\mu}\right)\mathcal O_i
 -C_i\gamma_{ij}^{\mathcal O}\mathcal O_j.
\end{align*}
把第二项的哑指标交换到与第一项相同的算符基，逐个比较 $\mathcal O_i$ 的系数即得正文\eqnrefs{eq:wilson-rg-dp11}。
\end{proof}

多个局域算符之间的尺度演化都具有这一矩阵结构；具体物理系统只改变算符基与反常维数矩阵。

复合算符的尺度演化确定以后，短距离乘积可以继续按同一局域算符基展开，从而把短程系数与长程矩阵元分离。

当两个重整化局域算符的类空间隔 $x$ 远小于强子、介质或外态的长程尺度时，局域性允许把短距离乘积按局域算符基作渐近展开。算符乘积展开（OPE）的结构为
\begin{equation}
 \boxed{
 \mathcal O_A(x)\mathcal O_B(0)
 \sim\sum_i C_{AB}^{\ \ i}(x,\mu)\mathcal O_i(0,\mu),
 \qquad x\to0.}
 \label{eq:ope-general-dp11}
\end{equation}
Wilson 系数 $C_{AB}^{\ i}$ 编码距离尺度 $|x|$ 附近的短程物理，在硬尺度足够高时可用微扰理论计算。外态矩阵元 $\langle H|\mathcal O_i|H\rangle$ 则承载长程动力学。

OPE 按算符维数与对称性排序。截断误差由下一批被省略算符相对于保留项的幂次决定，而不能只看某一个 Wilson 系数是否数值很小。

\eqnrefs{eq:ope-general-dp11} 的两边必须具有相同的重整化尺度依赖。复合算符的混合因此强制 Wilson 系数满足耦合的矩阵重整化群方程；选择 $\mu$ 接近短距离硬尺度可以把大对数从固定阶系数中转移到尺度演化中。具体的短程--长程分离问题只需把相应算符基和外态代入这一一般结构。

自由场给出一个最直接的检查。Wick 定理给
\begin{equation*}
 \varphi(x)\varphi(0)
 =\mathcal C(x)+:\!\varphi(x)\varphi(0)\!:.
\end{equation*}
对正规序部分在 $x=0$ 附近作 Taylor 展开，
\begin{equation*}
 :\!\varphi(x)\varphi(0)\!:
 =:\!\varphi^2(0)\!:
 +x^\mu:\!(\partial_\mu\varphi)\varphi\!:(0)+\cdots.
\end{equation*}
单位算符的 Wilson 系数就是短距离奇异的传播子 $\mathcal C(x)$，而 $:\!\varphi^2\!:$ 等局域算符承载较软的场信息。相互作用理论中的 OPE 保留这一职责分离，但系数和算符都必须经过重整化并允许混合。

零温连续频率积分在热平衡态中还会因虚时间边界条件离散化。热密度算符由
\begin{equation*}
 \hat\rho_{\beta_T}=\frac{\ee^{-\beta_T\hat H}}{Z(\beta_T)},
 \qquad Z(\beta_T)=\Tr \ee^{-\beta_T\hat H},
 \qquad \beta_T=(k_BT)^{-1}
\end{equation*}
定义。把实时间解析延拓到虚时间 $t=-\ii\tau$ 后，迹的循环性给出 Kubo--Martin--Schwinger 边界条件：玻色场在 $\tau\to\tau+\beta_T\hbar$ 下周期，费米场反周期，
\begin{align}
 \phi(\tau+\beta_T\hbar)&=+\phi(\tau),\\
 \psi(\tau+\beta_T\hbar)&=-\psi(\tau).
 \label{eq:matsubara-boundary-dp11}
\end{align}
因此虚时 Fourier 模式只能取离散的 Matsubara 频率
\begin{equation}
 \boxed{
 \begin{aligned}
 \omega_n^{(B)}&=\frac{2\pi n}{\beta_T\hbar},\\
 \omega_n^{(F)}&=\frac{(2n+1)\pi}{\beta_T\hbar}.
 \end{aligned}}
 \label{eq:matsubara-frequencies-dp11}
\end{equation}

\begin{proposition}[周期与反周期 Fourier 频率]
在欧氏时间长度为 $\beta\hbar$ 的区间上，函数满足相应周期或反周期边界条件。则有以下结论（原式 \eqnrefs{eq:matsubara-frequencies-dp11}）：
\begin{equation*}
 \boxed{
 \begin{aligned}
 \omega_n^{(B)}&=\frac{2\pi n}{\beta_T\hbar},\\
 \omega_n^{(F)}&=\frac{(2n+1)\pi}{\beta_T\hbar}.
 \end{aligned}}
\end{equation*}
\end{proposition}
\begin{proof}
玻色 Fourier 模写成 $\ee^{-\ii\omega_n\tau}$。周期条件给

\begin{equation*}
 \ee^{-\ii\omega_n\beta_T\hbar}=1,
\end{equation*}
因此
\begin{equation*}
 \omega_n^{(B)}=\frac{2\pi n}{\beta_T\hbar},
 \qquad n\in\mathbb Z.
\end{equation*}
费米场满足反周期条件，于是
\begin{equation*}
 \ee^{-\ii\omega_n\beta_T\hbar}=-1,
\end{equation*}
从而
\begin{equation*}
 \omega_n^{(F)}=\frac{(2n+1)\pi}{\beta_T\hbar}.
\end{equation*}
两者合起来就是正文\eqnrefs{eq:matsubara-frequencies-dp11}。在区间 $0\le\tau<\beta_T\hbar$ 上，这些模式满足
\begin{equation*}
 \frac1{\beta_T\hbar}\int_0^{\beta_T\hbar}\dd\tau\,
 \ee^{\ii(\omega_n-\omega_m)\tau}=\delta_{nm},
\end{equation*}
因此零温连续 Fourier 测度相应变为正文\eqnrefs{eq:matsubara-loop-replacement-dp11} 的离散求和。
\end{proof}

零温圈积分中的连续频率测度相应替换为
\begin{equation}
 \boxed{
 \int\frac{\dd\omega}{2\pi}
 \longrightarrow
 \frac1{\beta_T\hbar}\sum_n.}
 \label{eq:matsubara-loop-replacement-dp11}
\end{equation}
例如自由玻色子的虚时传播子分母为 $(\omega_n^{(B)})^2+\omega_{\mathbf k}^2$。完成频率求和后出现 Bose--Einstein 分布；统计因子由热边界条件与传播子谱极点共同决定。



多圈 $R$ 运算要求紫外反项保持局域；复合算符重整化把重合点奇异性组织成混合矩阵；算符乘积展开再把短距离奇异性编码进 Wilson 系数。三者描述的是同一局域量子场论在不同尺度组织方式下的同一个紫外结构。

